\section{Block averaging on the actual support}
\label{sec:block-averaging}

Let
\[
 \mathcal B=\{B_1,\ldots,B_m\}
\]
be a partition of \(\Omega\) into nonempty blocks.  For
\(\omega\in B_j\), define
\begin{equation}
 (P_{\mathcal B}A)_\omega
 =
 \overline A_j
 :=
 \frac{1}{|B_j|}
 \sum_{\nu\in B_j}A_\nu.
\label{eq:block-average}
\end{equation}

\begin{proposition}[Block averaging is a safe projection]
\label{prop:block-safe}
The operator \(P_{\mathcal B}\) is the orthogonal projection onto
the subspace of vectors that are constant on each block.  Moreover,
\[
 P_{\mathcal B}\one=\one.
\]
Consequently,
\begin{align}
 Z(P_{\mathcal B}A)&=Z(A),
\label{eq:block-zero-mode}\\
 D(P_{\mathcal B}A)
 &=
 \sum_{j=1}^m |B_j|\,|\overline A_j|^2,
\label{eq:block-energy}\\
 D(A)-D(P_{\mathcal B}A)
 &=
 \sum_{j=1}^m\sum_{\omega\in B_j}
 |A_\omega-\overline A_j|^2.
\label{eq:within-block-variance}
\end{align}
If \(P_{\mathcal B}A\ne0\), then
\[
 \Flat(A)\le\Flat(P_{\mathcal B}A),
\qquad
 \rho(A)\le\rho(P_{\mathcal B}A),
\]
where both statistics use the original frame size \(S\).
\end{proposition}

\begin{proof}
The formula \eqref{eq:block-average} is idempotent and its range is
the block-constant subspace.  For every block-constant \(C\),
\[
 \langle A-P_{\mathcal B}A,C\rangle
 =
 \sum_{j=1}^m
 \overline{C_j}
 \sum_{\omega\in B_j}
 (A_\omega-\overline A_j)
 =0.
\]
Hence the projection is orthogonal.  It clearly fixes \(\one\).
Equation \eqref{eq:block-zero-mode} follows by summing block means
with their multiplicities.  Equations \eqref{eq:block-energy} and
\eqref{eq:within-block-variance} are Pythagoras written block by
block.  The final conclusion is
\cref{thm:safe-projection-law}.
\end{proof}

\begin{corollary}[Exact flatness gap]
\label{cor:block-flatness-gap}
If \(Z(A)\ne0\) and \(P_{\mathcal B}A\ne0\), then
\begin{equation}
 \Flat(P_{\mathcal B}A)-\Flat(A)
 =
 \frac{|Z(A)|^2
 \bigl(D(A)-D(P_{\mathcal B}A)\bigr)}
 {D(P_{\mathcal B}A)D(A)}.
\label{eq:flatness-gap}
\end{equation}
Thus the inflation is exactly controlled by the discarded
within-block variance.
\end{corollary}

\begin{proof}
Subtract the two fractions with the common numerator
\(|Z(A)|^2\), and use \eqref{eq:within-block-variance}.
\end{proof}

\subsection{Nested coarse-graining}

Say that \(\mathcal B_{\rm fine}\) refines
\(\mathcal B_{\rm coarse}\) if every fine block is contained in a
coarse block.  Then
\[
 \Ran P_{\rm coarse}
 \subset
 \Ran P_{\rm fine}.
\]

\begin{theorem}[Monotone certificate ladder]
\label{thm:nested-monotonicity}
Let \(P_{\rm fine}\) and \(P_{\rm coarse}\) be the block projections
of nested partitions, with the fine partition refining the coarse
partition.  Then
\begin{equation}
 D(P_{\rm coarse}A)
 \le D(P_{\rm fine}A)
 \le D(A),
\label{eq:nested-energy}
\end{equation}
and all three vectors have the same zero mode.  On every nonzero
branch,
\begin{equation}
 \boxed{
 \Flat(A)
 \le\Flat(P_{\rm fine}A)
 \le\Flat(P_{\rm coarse}A),}
\label{eq:nested-flatness}
\end{equation}
with the analogous inequalities for \(\rho\), all computed on the
same frame \(\Omega\).
\end{theorem}

\begin{proof}
Because the ranges are nested,
\[
 P_{\rm coarse}P_{\rm fine}=P_{\rm coarse}.
\]
Apply norm contraction first to \(P_{\rm fine}A\), and then to
\(A\), to obtain \eqref{eq:nested-energy}.  Both projections fix
constants, hence preserve \(Z(A)\).  Dividing the common squared
zero mode by the ordered energies proves
\eqref{eq:nested-flatness}.
\end{proof}

\begin{remark}[Safe false negatives]
\label{rem:false-negatives}
A coarse block certificate may fail because its denominator has
lost too much within-block energy.  Refining the partition can only
improve the projected upper ratio.  Conversely, if even a coarse
level meets the desired upper bound, every finer level and the
physical coefficient meet it.  The ladder can therefore locate the
coarsest provably sufficient resolution without risking a false
positive.
\end{remark}

\subsection{Compression requires weights}

The range of \(P_{\mathcal B}\) has \(m\) degrees of freedom, but
replacing it by an unweighted vector in \(\C^m\) changes both the
zero mode and the Hilbert norm.  The correct compressed space is
\[
 H_{\mathcal B}
 =
 \C^m,\qquad
 \langle x,y\rangle_{\mathcal B}
 =
 \sum_{j=1}^m |B_j|x_j\overline{y_j}.
\]

\begin{proposition}[Weighted compressed isometry]
\label{prop:weighted-compression}
The map
\[
 C_{\mathcal B}:\Ran P_{\mathcal B}\to H_{\mathcal B},
\qquad
 C_{\mathcal B}(P_{\mathcal B}A)
 =
 (\overline A_1,\ldots,\overline A_m),
\]
is an isometry.  In compressed coordinates the physical statistics
are
\begin{align}
 Z_{\rm phys}
 &=
 \sum_{j=1}^m |B_j|\overline A_j,
\label{eq:weighted-zero}\\
 D_{\rm phys}
 &=
 \sum_{j=1}^m |B_j||\overline A_j|^2,
\label{eq:weighted-energy}\\
 S_{\rm phys}
 &=
 \sum_{j=1}^m |B_j|=S.
\label{eq:weighted-support}
\end{align}
These reproduce exactly the same-frame projected statistics.
\end{proposition}

\begin{proof}
Equations \eqref{eq:weighted-energy} and
\eqref{eq:weighted-support} give the norm identity, while
\eqref{eq:weighted-zero} is \eqref{eq:block-zero-mode}.
\end{proof}

\begin{example}[Unweighted compression manufactures cancellation]
\label{ex:unweighted-false-cancellation}
Let
\[
 \Omega=\{1,2,3\},\qquad
 B_1=\{1\},\quad B_2=\{2,3\},
\]
and take the already block-constant vector
\[
 A=(1,-1,-1).
\]
Then
\[
 Z_{\rm phys}=1-2=-1,\qquad
 D_{\rm phys}=1+2=3,\qquad
 \Flat_{\rm phys}=\frac13.
\]
The unweighted compressed vector \((1,-1)\), however, has apparent
sum \(0\) and apparent flatness \(0\).  No averaging error occurred:
the false cancellation was created solely by deleting the
multiplicities \((1,2)\).  Weighted compression restores the exact
physical values.
\end{example}

\begin{remark}[Atomic normalization rule]
\label{rem:atomic-normalization}
If literal atoms carry prescribed positive weights rather than unit
weights, the same theorem holds in the weighted Hilbert space and
block means must use those weights.  Changing from physical atomic
weights to unit block weights is a change of coefficient, not a
harmless coordinate compression.
\end{remark}
